% translated-from: en/chapters/04_paper_pipeline.tex@e4a959d4405d3a0ff9d71cc2cef14e592a31610d1bd8bee71993349cd5d0c271
% glossary:        5bca8d57a70f021e20c0af79aa06cd21013495c1a8246887ae889cd8a056f5fd
% model:           claude-opus-4-7

\section{論文生成パイプライン}\label{sec:paper}

\subsection{パイプライン概観}\label{sec:paper:overview}

\begin{figure}[t]
  \centering
  \resizebox{\columnwidth}{!}{% F06: paper-generation swim lane — paper-style three-lane workflow.
% Lanes are coloured stage containers; numbered step markers sit to the
% left of each station; the orange refine back-edge is on top.
% style.tikzstyles is loaded by shared/preamble.tex / standalone wrapper.
\begin{tikzpicture}[node distance=1.0cm and 1.6cm]

  % --- six stations laid out in a 3-row grid -------------------------
  \node[station]                                       (i1)  {\figlabel{F6.brainstorm}};
  \node[station, right=2.6cm of i1]                    (i2)  {\figlabel{F6.refine_idea}};

  \node[station emph, below=1.4cm of i1]               (e1)  {\figlabel{F6.bfts}};
  \node[station, right=2.6cm of e1]                    (e2)  {\figlabel{F6.replicate}};

  \node[station, below=1.4cm of e1]                    (w1)  {\figlabel{F6.draft}};
  \node[station emph, right=2.6cm of w1]               (w2)  {\figlabel{F6.review_refine}};

  % --- numbered step markers, perched on the upper-left corner of each
  % station (PaperBench-style badge), so they never overlap any arrow.
  \node[numbered step, above left=0.05cm and -0.2cm of i1.north west] {1};
  \node[numbered step, above left=0.05cm and -0.2cm of i2.north west] {2};
  \node[numbered step, above left=0.05cm and -0.2cm of e1.north west] {3};
  \node[numbered step, above left=0.05cm and -0.2cm of e2.north west] {4};
  \node[numbered step, above left=0.05cm and -0.2cm of w1.north west] {5};
  \node[numbered step, above left=0.05cm and -0.2cm of w2.north west] {6};

  % --- intra-lane forward arrows (thick, paper-style) ----------------
  \draw[big arrow] (i1) -- (i2);
  \draw[big arrow] (e1) -- (e2);
  \draw[big arrow] (w1) -- (w2);

  % --- inter-lane forward arrows: short straight verticals ----------
  \draw[big arrow] (i2.south) -- (e2.north);
  \draw[big arrow] (e2.south) -- (w2.north);

  % --- refine back-edge (orange, on top, label in middle) ----------
  \draw[big arrow, draw=ari-vermilion] (w2.south) to[bend left=35]
       node[edge label, fill=ari-orange!15, font=\sffamily\bfseries\footnotesize, midway]
       {\figlabel{F6.refine}}
       (w1.south);

  % --- lane backgrounds (in background layer) ----------------------
  \begin{scope}[on background layer]
    \node[stage container, fit=(i1)(i2), inner sep=10pt]   (laneI) {};
    \node[stage container, fit=(e1)(e2), inner sep=10pt]   (laneE) {};
    \node[stage container, fit=(w1)(w2), inner sep=10pt]   (laneW) {};
  \end{scope}
  % lane labels OUTSIDE on the right, after the lanes are drawn (top layer)
  \node[stage title, right=4pt of laneI.east, anchor=west] {\figlabel{F6.lane.idea}};
  \node[stage title, right=4pt of laneE.east, anchor=west] {\figlabel{F6.lane.experiment}};
  \node[stage title, right=4pt of laneW.east, anchor=west] {\figlabel{F6.lane.write}};
\end{tikzpicture}
}
  \caption{論文生成スイムレーン:着想、実験、執筆。各レーンが一
           つのフェーズに対応する;\texttt{review/refine} から
           \texttt{draft} へのバックエッジが唯一の非単調遷移であ
           る (\cref{sec:paper:loop})。}
  \label{fig:swim}
\end{figure}

論文生成パイプラインは ARI における二本目の活動の木である
(\cref{fig:swim})。探索が停止すると、生き残った系譜が
\term{paper スキル}に手渡されて草稿が生成され、続いて
\term{review スキル}が草稿を採点する。実装は
\code{ari-core/ari/pipeline/orchestrator.py:131} の
\texttt{run\_pipeline} 関数であり、同ファイル 68 行目の
\texttt{build\_scientific\_data} ヘルパが、paper スキルが必要と
する成果物群をノード単位で集約する。

一回の実行は永続化される 3 つの成果物を生む:\texttt{draft.tex}
(LaTeX 論文本体)、\texttt{review.json}(葉ごとのルブリックスコ
アと根拠)、そして \texttt{rubric.json}(本実行で使用したルブリ
ックインスタンス;ルブリックを動的生成する場合に有用)。

\subsection{ルブリック形式}\label{sec:paper:rubric}

\term{ルブリック} $\Rubric$ を、葉が組
$(\rubricitem{i},\rubricweight{i},\judgefn{i})$ を保持する有限木
としてモデル化する。各葉 $i$ は、カテゴリ記述子
$\rubricitem{i}$、非負の重み $\rubricweight{i}$
($\sum_i \rubricweight{i}=1$ を満たす)、そして段階的スコアを
返す\emph{判官関数} $\judgefn{i}:\paper\to[0,1]$ を持つ。論文ス
コアは凸結合
\begin{equation}\label{eq:paperscore}
  \paperscore{\paper} \;=\; \sum_i \rubricweight{i}\,\judgefn{i}(\paper).
\end{equation}
で与えられる。
この形式は PaperBench のルブリック木と共有されており、ARI は
これを vendoring された PaperBench パッケージを通じて直接再利
用する(\cref{sec:paper:paperbench}):葉型 \texttt{TaskNode} は
\texttt{paperbench.rubric.tasks} 由来、graded leaf 型
\texttt{GradedTaskNode} は \texttt{paperbench.judge.graded\_task\_node}
由来、重み付き集約エンベロープは
\code{ari-skill-paper-re/src/\_paperbench\_bridge.py:75} の
\texttt{aggregate\_graded\_tree} 由来である。

本番では二種類の判官族を用いる:LLM ベース判官は ARI 用に
PaperBench の \texttt{SimpleJudge}
(\code{ari-skill-paper-re/src/\_paperbench\_bridge.py:54})をラ
ップして実装する;決定的判官は二値または数値検証器を持つ項目
(例えば「再現可能性」「ドル絶対値で報告したコスト」)に使う。
両者は \texttt{Evaluator} プロトコル
(\code{ari-core/ari/protocols/evaluator.py:19})を実装する。

\subsection{review/refine ループ}\label{sec:paper:loop}

改稿サイクルはフィードバック $J(\paper_t)$ の下で
$\paper_t$ を $\paper_{t+1}$ へと反復的に書き換える:
\begin{equation}\label{eq:refine}
  \paper_{t+1} \;=\; \refine\!\bigl(\paper_t, J(\paper_t)\bigr).
\end{equation}

\term{収束基準}
$\paperscore{\paper_{t+1}}-\paperscore{\paper_t}<\epsilon$ が連
続二回成立した時、またはパイプラインあたり 3 回の上限に達した時
にサイクルは停止する。各分岐 — 収束 vs.\ 上限 — がどの程度の
頻度で発火するかは、凍結チェックポイントが揃い次第
\cref{chapter:eval} で報告する予定であり、本章では定量主張を行
わない。

改稿は意図的に軸ごとに非単調である:明確性の修正が数値軸を僅かに
劣化させ得る。これが \cref{fig:swim} のバックエッジの源であり、
収束基準は集約スコアにかかるのであって軸ごとではない。

\subsection{paper-re による PaperBench 統合}\label{sec:paper:paperbench}

\term{paper-re スキル}(\texttt{ari-skill-paper-re})は ARI の
PaperBench 互換複製ドライバである。統合は\emph{再実装ではない}:
PaperBench は \texttt{vendor/paperbench} に vendoring されて出
荷され、これにより ARI は upstream のタスク・ルブリック意味論を
逐語的に追跡する(配置要件は PaperBench の
\texttt{rubric.md \S{}8} に文書化されている)。3 つの統合点を述
べる。

\subsubsection*{1. ソルバ サブクラス: \texttt{AriPBSolver}}

\code{ari-skill-paper-re/src/\_replicator\_agent.py:269} は
PaperBench upstream の
\texttt{paperbench.solvers.basic.BasicAgentSolver} をサブクラス
化する。chz の immutable-config 規則により新フィールドが追加さ
れなくてもサブクラスにデコレータが必要となるが、ARI のオーバー
ライドは最小限である:

\begin{itemize}
  \item \texttt{\_get\_tools()}:submit 以外の全ツールを
        \texttt{\_BoundedOutputTool} でラップし、bash / python /
        file-read / search の出力を次の API 呼び出しへ流す前にサ
        イズ制限する。\texttt{submit} は名前で dispatch され、
        \texttt{execute()} が呼ばれることがないので、ラップされ
        ずにそのまま通過する。
  \item \texttt{\_run\_agent()}:
        \texttt{paperbench.solvers.utils.sanity\_check\_docker}
        を \texttt{\_bypass\_docker\_sanity\_check} で迂回し、
        upstream のハードコード
        \texttt{/home/\{submission,paper\}} パスを ARI のワーク
        スペース相対レイアウト
        (\texttt{LocalComputer} cwd = \texttt{repro\_sandbox/}、
        論文を \texttt{paper/} に、submission を \texttt{submission/}
        に置く)へ書き換える。プロンプト内容はその他の点で
        upstream の \texttt{get\_instructions} /
        \texttt{get\_system\_message} から逐語的に取得され、
        ARI は upstream の複製タスク枠組み(7 日予算、「すべての
        結果を複製するまで停止しない」等)に整合する。
\end{itemize}

\texttt{AriPBSolver} はまた upstream の
\texttt{check\_for\_existing\_run} を保持しており、部分的に完了
した実行を冪等に再開できる。

\subsubsection*{2. ルブリック駆動の成果物リスト}

vendor PaperBench は ARI の論文ごとのルブリック木を知らないた
め、\texttt{AriPBSolver.\_run\_agent} は
\texttt{LocalPBTask.rubric\_expected\_artifacts} が空でない場合、
ユーザ指示に \texttt{EXPECTED\_ARTIFACTS} ブロックを追記する。
追記されるブロックは、成功した \texttt{reproduce.sh} が産出すべ
きパスを列挙する;これがないと、grader の葉主張はどのファイル
を検証すべきかのヒントを欠く。ルブリック木自体は
\code{ari-skill-paper-re/src/\_paperbench\_bridge.py:63}
(\texttt{task\_node\_from\_dict})で構築される。

\subsubsection*{3. ドライバ: \texttt{run\_replicator\_agent}}

\code{ari-skill-paper-re/src/\_replicator\_agent.py:365} が非同
期エントリポイントである。標準の \texttt{build\_reproduce\_sh}
エンベロープ
$\{$\texttt{populated},
\texttt{output\_dir}, \texttt{files},
\texttt{expected\_artifacts}, \texttt{max\_runtime\_sec},
\texttt{model}, \texttt{prompt\_sha256}, \texttt{notes},
\texttt{warnings}$\}$ を返す。\texttt{output\_dir} はエージェ
ントのワークスペースであると同時に、ARI が rollout 後に
\texttt{reproduce.sh} を見つけることを期待する場所でもある;
upstream のソルバは \texttt{agent.log} 等の bookkeeping を
\texttt{run\_dir} に書き込むため、これをワークスペース自身に向
けることで成果物が \texttt{reproduce.sh} と隣接する。

既定の completer は \texttt{OpenAIResponsesTurnCompleterConfig}
でモデルは \texttt{gpt-5-mini}(環境変数
\texttt{ARI\_MODEL\_REPLICATOR} で上書き可能);LiteLLM ルート
の代替は \texttt{\_litellm\_completer.py} のヘルパ経由で配線さ
れる。既定の時間上限は実行あたり 12 時間;\texttt{sandbox\_kind}
は \texttt{auto} で、ローカル Apptainer イメージが供給されてい
れば fallback する。

\subsubsection*{4. 採点と集約}

複製エージェントが完了すると、採点は
\code{ari-skill-paper-re/src/\_paperbench\_bridge.py} の
\texttt{judge\_submission} に委譲される。この関数は ARI の
$($\texttt{paper\_md}, \texttt{submission\_dir},
\texttt{reproduce\_log}, \texttt{judge\_model}$)$ タプルを
upstream の \texttt{SimpleJudge} に適合させ、submission ごとに
\texttt{GradedTaskNode} 木を生成する。実行間の集約は次の二つの
ヘルパで処理される:

\begin{itemize}
  \item \texttt{aggregate\_graded\_tree}(line~75):一つの
        \texttt{GradedTaskNode} をルブリック重み
        $\rubricweight{i}$ 経由で重み付きスカラーへ畳み込み、
        カテゴリ別内訳を併せて返す。
  \item \texttt{average\_graded\_runs}(line~91 周辺):$n$ 個
        の \texttt{GradedTaskNode} 木に対する実行ごとの平均を取
        り、ノイズ低減のためエージェントを再走させた場合でも単
        一の複製スコアが報告される。
\end{itemize}

二段階設計 — 葉採点には vendor の \texttt{SimpleJudge}、形状と
集約には ARI のアダプタ — により、PaperBench upstream のアップ
グレードは vendor swap となり(ARI コード変更不要)、ARI は葉採
点を単一数値へ結合する仕方を完全にコントロールする。同じエンベ
ロープが葉が二値のとき $\judgefn{i}\in\{0,1\}$ として
\cref{eq:paperscore} に流れる。

\subsection{故障モードとガードレール}\label{sec:paper:failures}

頻出する故障モードを 3 つ命名する:

\begin{itemize}
  \item \textbf{LLM のみの裏付け}:論文が、系譜中のいずれのノー
        ドも裏付けない主張をする。緩和策は静的検査:数値主張ま
        たは因果主張は、ノード成果物への
        \texttt{\textbackslash code\{...\}} 参照経由で必ず辿れ
        ねばならない。
  \item \textbf{引用ハルシネーション}:論文が存在しない論文を引
        用する。緩和策は構造的:LLM が引用キーを発明することを
        禁じ、すべての bib エントリは Semantic Scholar に対して
        検証されなければならない (\cref{sec:related})。
  \item \textbf{改稿ドリフト}:スコアは上昇するが論文が一貫性を
        失う。緩和策はルブリックの軸ごと下限:いずれかの軸が
        $0.5$ を下回れば改稿は中止される。
\end{itemize}

草稿時のソースファイル選択は \code{node\_selection.py:204}
(\texttt{select\_source\_files\_for\_publication})が担う;主張
が頼ることのできる系譜ノードを正確に列挙することで、LLM のみの
裏付け型の故障を防ぐ。コンパニオンの
\texttt{load\_selected\_sources}(line~258)が
\texttt{(node\_id, rel\_path)} 各組のバイトを読み、オプションの
\texttt{size\_budget} を尊重するので、論文草稿が context を超え
ることは決してない。

\subsection{複製器のためのツール出力境界}\label{sec:paper:bounded-tool}

複製エージェントは単一の bash セッション内で何時間も過ごし得
る;ツール出力(長い stdout、大規模ディレクトリリスト)は放置
すれば後続の各プロンプトを支配する。ARI は
\texttt{\_BoundedOutputTool}
(\code{ari-skill-paper-re/src/\_replicator\_agent.py:243})でこ
れを境界し、同ツールが内部ツールの \texttt{execute()} をラッ
プして出力を \texttt{\_truncate\_tool\_output}(line~218)で構
成可能なバイト上限へ後処理する。これが PaperBench upstream の
自由形式トランスクリプトと ARI の境界付き再生可能なものとの違
いであり、複数時間の rollout でもプロンプト予算を予測可能に保
つ理由である。
